% Bestiary template for D&D 5e creatures with grouping support
% Uses dndbook document class with structured creature organization

% Base LaTeX template for D&D 5e documents using DND-5e-LaTeX-Template
% This template provides the foundation for all D&D-style documents with structured content

\documentclass[letterpaper, 11pt, bg=full, twocolumn, stats=modern]{dndbook}

% DND-5e-LaTeX-Template provides most styling automatically
% Only include additional packages if needed

% Additional packages for enhanced functionality
\usepackage{graphicx}
\usepackage{wrapfig}
\usepackage{hyperref}
\usepackage{bookmark}
\usepackage{makeidx}

% Additional packages for bestiary layout
\usepackage{multicol}
\usepackage{enumitem}
\usepackage{tcolorbox}
\usepackage{titlesec}

% Enhanced list environments for creature details
\tcbuselibrary{breakable}

% Float management for large bestiaries
\usepackage{afterpage}
\usepackage{placeins}
\usepackage{stfloats} % enable bottom placement for double-column floats

\title{Creature Bestiary}
\date{\today}

% Enable automatic numbering for chapters only, not sections
\setcounter{secnumdepth}{0}

% Disable text justification but preserve paragraph indentation
\raggedright
\raggedbottom
% Restore paragraph indentation after \raggedright resets it
\setlength{\parindent}{0.8em}

% Hyperlink configuration
\hypersetup{
    colorlinks=true,
    linkcolor=black,
    urlcolor=black,
    citecolor=black,
    filecolor=black,
    bookmarks=true,
    bookmarksopen=true,
    bookmarksnumbered=true,
    pdftitle={Creature Bestiary},
    pdfsubject={D&D 5e Content},
    pdfkeywords={D&D, 5e, RPG, tabletop, }
}

% Enhanced section styling for bestiary
% DND-5e-LaTeX-Template provides all styling automatically

% Configure for bestiary structure
\setcounter{secnumdepth}{0}  % No section numbers for cleaner look
\setcounter{tocdepth}{2}     % Include subsections in table of contents

% Studiorum Image Helpers — simple, explicit commands
% Load required packages locally so templates don't need to manage them.
\RequirePackage{graphicx}
\RequirePackage{xparse}
\RequirePackage{caption}

% Inline (non-floating) image block
% Usage: \StudiorumImageInline[<width fraction>]{file}{caption}[label]
\NewDocumentCommand{\StudiorumImageInline}{ s o m m o }{%%
  \par\medskip
  \noindent\centering
  \IfNoValueTF{#2}{%%
    \includegraphics[width=\linewidth,height=0.45\textheight,keepaspectratio]{#3}%%
  }{%%
    \includegraphics[width=#2\linewidth,height=0.45\textheight,keepaspectratio]{#3}%%
  }%%
  \IfBooleanF{#1}{%%
    \par\smallskip
    \captionof{figure}{#4}%%
    \IfNoValueF{#5}{\label{#5}}%%
  }%%
  \par\medskip
}

% Standard floating figure
% Usage: \StudiorumImageFloat[<width fraction>]{file}{caption}[label]
\NewDocumentCommand{\StudiorumImageFloat}{ s o m m o }{%%
  \begin{figure}[htbp]
    \centering
    \IfNoValueTF{#2}{%%
      \includegraphics[width=\linewidth,height=0.45\textheight,keepaspectratio]{#3}%%
    }{%%
      \includegraphics[width=#2\linewidth,height=0.45\textheight,keepaspectratio]{#3}%%
    }%%
    \IfBooleanF{#1}{%%
      \caption{#4}%%
      \IfNoValueF{#5}{\label{#5}}%%
    }%%
  \end{figure}%%
}

% Wide (two-column spanning when applicable)
% Usage: \StudiorumImageWide[<width fraction>]{file}{caption}[label]
\NewDocumentCommand{\StudiorumImageWide}{ s o m m o }{%%
  \ifdim\columnwidth=\textwidth
    \begin{figure}[htbp]
      \centering
      \IfNoValueTF{#2}{%%
        \includegraphics[width=\linewidth,height=0.45\textheight,keepaspectratio]{#3}%%
      }{%%
        \includegraphics[width=#2\linewidth,height=0.45\textheight,keepaspectratio]{#3}%%
      }%%
      \IfBooleanF{#1}{%%
        \caption{#4}%%
        \IfNoValueF{#5}{\label{#5}}%%
      }%%
    \end{figure}%%
  \else
    \begin{figure*}[tp]
      \centering
      \IfNoValueTF{#2}{%%
        \includegraphics[width=\textwidth,height=0.45\textheight,keepaspectratio]{#3}%%
      }{%%
        \includegraphics[width=#2\textwidth,height=0.45\textheight,keepaspectratio]{#3}%%
      }%%
      \IfBooleanF{#1}{%%
        \caption{#4}%%
        \IfNoValueF{#5}{\label{#5}}%%
      }%%
    \end{figure*}%%
  \fi
}

% Full-page (dedicated page; fits within text block)
% Usage: \StudiorumImageFullpage[<width fraction>]{file}{caption}[label]
\NewDocumentCommand{\StudiorumImageFullpage}{ s o m m o }{%%
  \clearpage
  \begin{figure}[p]
    \centering
    \IfNoValueTF{#2}{%%
      \includegraphics[width=\textwidth,height=\textheight,keepaspectratio]{#3}%%
    }{%%
      \includegraphics[width=#2\textwidth,height=\textheight,keepaspectratio]{#3}%%
    }%%
    \IfBooleanF{#1}{%%
      \caption{#4}%%
      \IfNoValueF{#5}{\label{#5}}%%
    }%%
  \end{figure}
  \clearpage
}

% Full-page bleed (fills the physical page; may distort aspect ratio)
% Usage: \StudiorumImageFullpageBleed{file}{caption}[label]
\NewDocumentCommand{\StudiorumImageFullpageBleed}{ s o m m o }{%%
  \clearpage
  \begin{figure}[p]
    \centering
    \includegraphics[width=\paperwidth,height=\paperheight]{#3}%%
    \IfBooleanF{#1}{%%
      \caption{#4}%%
      \IfNoValueF{#5}{\label{#5}}%%
    }%%
  \end{figure}
  \clearpage
}

% No-caption convenience aliases
\NewDocumentCommand{\StudiorumImageInlineNoCaption}{ o m }{\StudiorumImageInline*[#1]{#2}{}}
\NewDocumentCommand{\StudiorumImageFloatNoCaption}{ o m }{\StudiorumImageFloat*[#1]{#2}{}}
\NewDocumentCommand{\StudiorumImageWideNoCaption}{ o m }{\StudiorumImageWide*[#1]{#2}{}}
\NewDocumentCommand{\StudiorumImageFullpageNoCaption}{ o m }{\StudiorumImageFullpage*[#1]{#2}{}}
\NewDocumentCommand{\StudiorumImageFullpageBleedNoCaption}{ m }{\StudiorumImageFullpageBleed*{#1}{}}

\begin{document}

\tableofcontents
\clearpage

\begin{DndComment}{About This Bestiary}
This bestiary contains 3 creatures from SRD.

\textbf{Challenge Ratings:} CR 1/4: 1, CR 17: 1, CR 21: 1

\textbf{Creature Types:} Dragon: 1, Humanoid (Goblinoid): 1, Undead: 1
\end{DndComment}

\vspace{1em}

\section{CR 1/8-1/4}

\vspace{0.5em}

\addcontentsline{toc}{subsection}{Goblin}

\begin{DndMonster}[float=tp]{Goblin}
\noindent \DndMonsterType{Small humanoid (goblinoid), neutral evil}
\DndMonsterBasics[
    armor-class = {15 (\textit{leather armor}, \textit{shield})},
    hit-points = {7 (2d6)},
    speed = {30 ft.},
    initiative = +2
]
\DndMonsterAbilityScores[
    str = 8,
    dex = 14,
    con = 10,
    int = 10,
    wis = 8,
    cha = 8,
]

\DndMonsterDetails[
    skills = {Stealth +6},
    senses = {darkvision 60 ft.},
    languages = {Common, Goblin},
    challenge = {1/4 (50 XP)}
]

\DndMonsterSection{Traits}
\DndMonsterAction{Nimble Escape}
The goblin can take the Disengage or Hide action as a bonus action on each of its turns.

\DndMonsterSection{Actions}
\DndMonsterAction{Scimitar}
Melee Weapon Attack: +4 to hit, reach 5 ft., one target. Hit: 5 (1d6 + 2) slashing damage.

\DndMonsterAction{Shortbow}
Ranged Weapon Attack: +4 to hit, range 80/320 ft., one target. Hit: 5 (1d6 + 2) piercing damage.

% Add source abbreviation at the bottom
\vspace{2pt}
\noindent\hfill{\scriptsize\textcolor{gray}{SRD}}

\end{DndMonster}

\vspace{2em}

\FloatBarrier

\section{CR 17-20}

\vspace{0.5em}

\addcontentsline{toc}{subsection}{Adult Red Dragon}

\begin{DndMonster}[float*=tp,width=\textwidth + 8pt, toptitle=0pt, top=1pt]{Adult Red Dragon}
\vspace*{-1.8\baselineskip}
\begin{multicols}{2}
\raggedcolumns
\begingroup
\setlength{\premulticols}{0pt}
\setlength{\postmulticols}{0pt}
\setlength{\multicolsep}{0pt}
\setlength{\parskip}{0pt}
\noindent \DndMonsterType{Huge dragon, chaotic evil}
\DndMonsterBasics[
    armor-class = {19 (natural armor)},
    hit-points = {256 (19d12 + 133)},
    speed = {40 ft., fly 80 ft., climb 40 ft.},
    initiative = +0
]
\DndMonsterAbilityScores[
    str = 27,
    dex = 10,
    con = 25,
    int = 16,
    wis = 13,
    cha = 21,
    dex save = +6,
    con save = +13,
    wis save = +7,
    cha save = +11
]

\DndMonsterDetails[
    skills = {Perception +13, Stealth +6},
    damage-immunities = {fire},
    senses = {blindsight 60 ft., darkvision 120 ft.},
    languages = {Common, Draconic},
    challenge = {17 (18,000 XP)}
]

\DndMonsterSection{Traits}
\DndMonsterAction{Legendary Resistance (3/Day)}
If the dragon fails a saving throw, it can choose to succeed instead.

\DndMonsterSection{Actions}
\DndMonsterAction{Multiattack}
The dragon can use its Frightful Presence. It then makes three attacks: one with its bite and two with its claws.

\DndMonsterAction{Bite}
Melee Weapon Attack: +14 to hit, reach 10 ft., one target. Hit: 19 (2d10 + 8) piercing damage plus 7 (2d6) fire damage.

\DndMonsterAction{Claw}
Melee Weapon Attack: +14 to hit, reach 5 ft., one target. Hit: 15 (2d6 + 8) slashing damage.

\DndMonsterAction{Tail}
Melee Weapon Attack: +14 to hit, reach 15 ft., one target. Hit: 17 (2d8 + 8) bludgeoning damage.

\DndMonsterAction{Frightful Presence}
Each creature of the dragon's choice that is within 120 feet of the dragon and aware of it must succeed on a DC 19 Wisdom saving throw or become frightened for 1 minute. A creature can repeat the saving throw at the end of each of its turns, ending the effect on itself on a success. If a creature's saving throw is successful or the effect ends for it, the creature is immune to the dragon's Frightful Presence for the next 24 hours.

\DndMonsterAction{Fire Breath (Recharge 5--6)}
The dragon exhales fire in a 60-foot cone. Each creature in that area must make a DC 21 Dexterity saving throw, taking 63 (18d6) fire damage on a failed save, or half as much damage on a successful one.

\DndMonsterSection{Legendary Actions}
The Adult Red Dragon can take 3 legendary actions, choosing from the options below. Only one legendary action can be used at a time and only at the end of another creature's turn. The Adult Red Dragon regains spent legendary actions at the start of its turn.

\DndMonsterAction{Detect}
The dragon makes a Wisdom (\textit{Perception}) check.

\DndMonsterAction{Tail Attack}
The dragon makes a tail attack.

\DndMonsterAction{Wing Attack (Costs 2 Actions)}
The dragon beats its wings. Each creature within 10 feet of the dragon must succeed on a DC 22 Dexterity saving throw or take 15 (2d6 + 8) bludgeoning damage and be knocked prone. The dragon can then fly up to half its flying speed.

% Add source abbreviation at the bottom
\vspace{2pt}
\noindent\hfill{\scriptsize\textcolor{gray}{SRD}}

\endgroup
\end{multicols}
\end{DndMonster}

\vspace{2em}

\FloatBarrier

\section{CR 21+}

\vspace{0.5em}

\addcontentsline{toc}{subsection}{Lich}

\begin{DndMonster}[float*=tp,width=\textwidth + 8pt, toptitle=0pt, top=1pt]{Lich}
\vspace*{-1.8\baselineskip}
\begin{multicols}{2}
\raggedcolumns
\begingroup
\setlength{\premulticols}{0pt}
\setlength{\postmulticols}{0pt}
\setlength{\multicolsep}{0pt}
\setlength{\parskip}{0pt}
\noindent \DndMonsterType{Medium undead, any evil alignment}
\DndMonsterBasics[
    armor-class = {17 (natural armor)},
    hit-points = {135 (18d8 + 54)},
    speed = {30 ft.},
    initiative = +3
]
\DndMonsterAbilityScores[
    str = 11,
    dex = 16,
    con = 16,
    int = 20,
    wis = 14,
    cha = 16,
    con save = +10,
    int save = +12,
    wis save = +9,
]

\DndMonsterDetails[
    skills = {Arcana +19, History +12, Insight +9, Perception +9},
    damage-resistances = {cold; lightning; necrotic},
    damage-immunities = {poison, bludgeoning, piercing, slashing},
    condition-immunities = {charmed, exhaustion, frightened, paralyzed, poisoned},
    senses = {truesight 120 ft.},
    languages = {Common plus up to five other languages},
    challenge = {21 (33,000 XP) or 22 (41,000 XP) when encountered in lair}
]

\DndMonsterSection{Traits}
\DndMonsterAction{Spellcasting}
The lich is an 18th-level spellcaster. Its spellcasting ability is Intelligence (spell save DC 20, +12 to hit with spell attacks). The lich has the following wizard spells prepared:
\textbf{Cantrips (at will):} \textit{mage hand}, \textit{prestidigitation}, \textit{ray of frost}
\textbf{1st level (4 slots):} \textit{detect magic}, \textit{magic missile}, \textit{shield}, \textit{thunderwave}
\textbf{2nd level (3 slots):} \textit{detect thoughts}, \textit{invisibility}, \textit{Melf's acid arrow}, \textit{mirror image}
\textbf{3rd level (3 slots):} \textit{animate dead}, \textit{counterspell}, \textit{dispel magic}, \textit{fireball}
\textbf{4th level (3 slots):} \textit{blight}, \textit{dimension door}
\textbf{5th level (3 slots):} \textit{cloudkill}, \textit{scrying}
\textbf{6th level (1 slots):} \textit{disintegrate}, \textit{globe of invulnerability}
\textbf{7th level (1 slots):} \textit{finger of death}, \textit{plane shift}
\textbf{8th level (1 slots):} \textit{dominate monster}, \textit{power word stun}
\textbf{9th level (1 slots):} \textit{power word kill}

\DndMonsterAction{Legendary Resistance (3/Day)}
If the lich fails a saving throw, it can choose to succeed instead.

\DndMonsterAction{Rejuvenation}
If it has a phylactery, a destroyed lich gains a new body in 1d10 days, regaining all its hit points and becoming active again. The new body appears within 5 feet of the phylactery.

\DndMonsterAction{Turn Resistance}
The lich has advantage on saving throws against any effect that turns undead.

\DndMonsterSection{Actions}
\DndMonsterAction{Paralyzing Touch}
Melee Spell Attack: +12 to hit, reach 5 ft., one creature. Hit: 10 (3d6) cold damage. The target must succeed on a DC 18 Constitution saving throw or be paralyzed for 1 minute. The target can repeat the saving throw at the end of each of its turns, ending the effect on itself on a success.

\DndMonsterSection{Legendary Actions}
The Lich can take 3 legendary actions, choosing from the options below. Only one legendary action can be used at a time and only at the end of another creature's turn. The Lich regains spent legendary actions at the start of its turn.

\DndMonsterAction{Cantrip}
The lich casts a cantrip.

\DndMonsterAction{Paralyzing Touch (Costs 2 Actions)}
The lich uses its Paralyzing Touch.

\DndMonsterAction{Frightening Gaze (Costs 2 Actions)}
The lich fixes its gaze on one creature it can see within 10 feet of it. The target must succeed on a DC 18 Wisdom saving throw against this magic or become frightened for 1 minute. The frightened target can repeat the saving throw at the end of each of its turns, ending the effect on itself on a success. If a target's saving throw is successful or the effect ends for it, the target is immune to the lich's gaze for the next 24 hours.

\DndMonsterAction{Disrupt Life (Costs 3 Actions)}
Each non-undead creature within 20 feet of the lich must make a DC 18 Constitution saving throw against this magic, taking 21 (6d6) necrotic damage on a failed save, or half as much damage on a successful one.

% Add source abbreviation at the bottom
\vspace{2pt}
\noindent\hfill{\scriptsize\textcolor{gray}{SRD}}

\endgroup
\end{multicols}
\end{DndMonster}

\FloatBarrier

\end{document}